% Page 100

\seite{100}

\abschnitt{1931}
\pstart
\margmark{15. Januar}%
Der Obstbauberater Blaes aus Malberg hielt dieser\ob
Tage eine Gründungsversammlung zu einem Obst\ohb
Landesbauverein abzuhalten. Angenehm erfreulich\ob
war der Lehrer\unsicher{} in. Vorsitzender des Vereins ward. Die\ob
Gründung des Vereins beschlossen. Der gleichzeitig ge\ohb
wählte Vorstand setzt sich wie folgt zusammen:\ob
Hillen Conrad, Bostpfeifer. Weinandt, Geissler,\ob
Bigelschiefer. Rothe Gerland, Gassenack und\ob
Klassen Wilhelm, Eiffgen.

\margmark{Februar}%
Am vergangenen Samstag hielt die Gemein\ohb
de eine dreitägige Holzversteigerung ab. Die über\ob
Erwarten stark besucht war. Es kam Nutz- und\ob
Brennholz zum Ausgebot. Das Ergebnis war ein\ob
befriedigendes. Die erzielten Durchschnittspreise\ob
waren folgende: Fichtenbrennholz 22,50 M.\ pro Festen\ohb
meter, Eichenstammholz 25,00 pro Festm., Birkenbrennholz 3,50 M.\ob
pro Rm., Fichtenknüppel 5,00 M.\ und Eichen-Knüppel 4 M.\ je\ob
pro Raummeter.

In der daran anschließend fand auf Vorschlag des\ob
Bürgermeisters eine Nachtragssitzung Gemeinderatssit\ohb
zung statt würde. Die dreitägige Gemeinderatssitzung, bei Voll\ohb
besetzter Klassen Gerhard anwesend mit Stimmen\ohb
mehrheit einstimmig den Beschluss gefasst: Die\ob
Waffe des Weltkrieges soll mit Schirmen bekannt.

Am selbigen Sonntag hielt der Krieger\ohb
verein Sefferweich, der Pfarre Sefferweich, im\ob
Vereinslokal seine ordentliche Gründungsversamm\ohb
lung ab. Da die am 18.\ Januar angesetzte erste Ver\ohb
günstigung Sitzung die Lage der Beschluss\unsicher...
\pend
